% Chapter 1 — Introduction

Turbulent reacting flows inside gas-turbine combustors involve a tight, multi-scale coupling between fluid mechanics, chemical kinetics, and heat release. Understanding and monitoring the spatial distribution of temperature inside these chambers is central to combustor design and operation: the temperature field governs flame stabilization, local quenching, wall thermal loading, thermoacoustic instabilities, and pollutant formation pathways. Yet in experimental test rigs, direct two-dimensional measurement of the instantaneous temperature field remains substantially harder and more restrictive than measuring velocity. This thesis investigates how accurately a neural reconstruction model can infer instantaneous two-dimensional temperature fields from spatially incomplete velocity observations---and optional sparse thermal measurements---across varying fuel blends and equivalence ratios in a swirl-stabilized ammonia--methane combustor.

\section{Industrial and Energy Context: Ammonia--Methane Co-Firing in Swirl Combustors}
\label{sec:intro_context}

Unlike reciprocating internal combustion engines, where combustion occurs intermittently inside a closed cylinder bounded by a moving piston, aviation and stationary power-generation gas turbines operate in a continuous open-flow regime. Air compressed by the axial compressor enters the combustion chamber continuously, mixes with fuel, and burns in a self-sustained flame before expanding through the turbine stages. Because bulk axial flow velocities inside the combustor far exceed the laminar or turbulent burning velocity of hydrocarbon and carbon-free fuels, a free jet flame would immediately blow out without an aerodynamic anchoring mechanism.

Modern gas-turbine combustors achieve flame anchoring through \emph{swirl stabilization}~\cite{syred1974combustion, candel2014dynamics}. By imparting a tangential velocity component $u_\theta$ to the incoming reactant stream via swirl vanes, centrifugal forces establish a radial pressure gradient ($\partial p / \partial r \approx \rho u_\theta^2 / r$) that depresses static pressure along the burner axis. As the swirling annular jet expands into the wider combustion chamber, the decay of tangential velocity causes this low-pressure core to recover downstream; once the swirl number exceeds a critical threshold ($S_w \gtrsim 0.6$), the resulting adverse axial pressure gradient ($\partial p / \partial x > 0$) overcomes the centerline axial momentum and triggers \emph{vortex breakdown}. This aerodynamic instability creates a toroidally recirculating low-velocity region along the axis known as the \emph{Inner Recirculation Zone} (IRZ), complemented by corner \emph{Outer Recirculation Zones} (ORZ) bounded by the combustor dump plane and liner walls. The IRZ continuously transports hot burnt products ($T \approx 1900\text{--}2280\text{ K}$) and active radical species backward toward the injector exit, providing a continuous thermal and chemical reservoir that ignites the incoming cold unburnt reactant jets ($T \approx 300\text{--}400\text{ K}$) along the sheared mixing layers.

In the transition toward low-carbon power generation and sustainable propulsion, ammonia ($\text{NH}_3$) has emerged as a prominent carbon-free hydrogen carrier~\cite{valera2018ammonia, kobayashi2019science}. Unlike pure hydrogen ($\text{H}_2$), which requires cryogenic storage at $-253\,^\circ\text{C}$ ($20\text{ K}$) or compression above $350\text{--}700\text{ bar}$, ammonia liquefies at modest pressure ($\sim 8.6\text{ bar}$ at $20\,^\circ\text{C}$) and benefits from an established global synthesis and distribution infrastructure. However, burning pure ammonia directly in a gas-turbine combustor poses severe combustion challenges:
\begin{itemize}
    \item \textbf{Low reactivity and laminar burning velocity:} The unstretched laminar burning velocity of stoichiometric $\text{NH}_3/\text{air}$ at atmospheric pressure ($\sim 0.065\text{--}0.07\text{ m/s}$) is roughly one-fifth that of methane ($\text{CH}_4/\text{air}$, $\sim 0.37\text{ m/s}$)~\cite{hayakawa2015laminar, okafor2018experimental}, accompanied by a higher minimum ignition energy and a narrower flammability range ($\phi \approx 0.63\text{--}1.40$ vs.\ $0.50\text{--}1.68$ for $\text{CH}_4$). Consequently, ammonia-rich flames are highly susceptible to lean blow-off and local extinction under high turbulent strain rates.
    \item \textbf{Reduced flame temperature and slower heat release:} At atmospheric stoichiometric conditions ($T_u = 300\text{ K}$), replacing methane with pure ammonia lowers the equilibrium adiabatic flame temperature by $\sim 150\text{ K}$ (from $\approx 2225\text{ K}$ to $\approx 2070\text{ K}$). In a strained turbulent swirl combustor, finite-rate chemical kinetics and slower heat release amplify this drop, reducing domain-mean temperatures by $185\text{--}225\text{ K}$ (and conditional temperatures at matched flow speed by $200\text{--}360\text{ K}$) as the ammonia fraction increases from $0\%$ to $80\%$, while lengthening the ignition delay and the penetration depth of the cold reactant jets.
    \item \textbf{Fuel-bound nitrogen emissions:} Unlike pure methane, where $\text{NO}_x$ forms predominantly via the high-temperature thermal (Zeldovich) pathway above $1800\text{ K}$, ammonia oxidation is governed by fuel-$\text{NO}$ kinetics ($\text{NH}_3 \rightarrow \text{NH}_2 \rightarrow \text{NH} \rightarrow \text{HNO} \rightarrow \text{NO}$) that depend sensitively on the local equivalence ratio $\phi$ and temperature distribution~\cite{kobayashi2019science}. Lean and stoichiometric operation can produce high $\text{NO}$ emissions, local flame quenching below $1400\text{ K}$ promotes nitrous oxide ($\text{N}_2\text{O}$, a potent greenhouse gas) formation via $\text{NH} + \text{NO}$ pathways, and fuel-rich operation ($\phi > 1.1$) promotes $\text{NO}$ reduction via $\text{NH}_i$ radicals at the cost of unburnt $\text{NH}_3$ slip and endothermic $\text{H}_2$ dissociation.
\end{itemize}

To bridge the gap between conventional natural gas operation and carbon-free firing, industrial turbine developers (such as GE Vernova, Mitsubishi Power, and Ansaldo Energia) actively test \emph{ammonia--methane ($\text{NH}_3/\text{CH}_4$) co-firing} across wide ranges of $\text{NH}_3$ blending ratios ($0\%$ to $80\%+$) and equivalence ratios $\phi$. Characterizing how the flame structure, recirculation zone temperature, and shear-layer fluctuations respond across this two-dimensional operating space $(\text{NH}_3\%, \phi)$ is therefore a primary objective of laboratory combustor testing.

\section{The Experimental Diagnostic Bottleneck in Reacting Flows}
\label{sec:intro_diagnostics}

When an optically accessible swirl combustor is tested in a propulsion or energy laboratory, researchers and engineers seek simultaneous knowledge of the velocity field $\mathbf{u}(\mathbf{x}, t)$ and the temperature field $T(\mathbf{x}, t)$ across a two-dimensional diametral plane. In practice, however, experimental diagnostics face a sharp asymmetry between kinematic and thermal measurements:

\begin{enumerate}
    \item \textbf{Velocity is accessible in 2D via Particle Image Velocimetry (PIV), but often incomplete:} By seeding the flow with refractory sub-micron ceramic particles (such as $\text{TiO}_2$, $\text{Al}_2\text{O}_3$, or $\text{ZrO}_2$) and illuminating a diametral plane with a high-speed double-pulsed laser sheet ($4\text{--}10\text{ kHz}$), planar PIV (2-component: radial and axial velocity $v_i, v_j$) or stereoscopic PIV (3-component, adding out-of-plane swirl $v_k$) yields instantaneous velocity vectors across the measurement plane~\cite{raffel2018particle}. Even in laboratory PIV, however, valid velocity vectors are rarely available uniformly across $100\%$ of the grid: strong gas expansion across the flame front ($\rho_u/\rho_b \approx 6\text{--}7$) dilutes tracer seeding density in the hot burnt products, while thermophoretic drift across steep thermal gradients, refractive-index beam steering, laser flare near quartz walls, or restricted optical access windows frequently leave only a sparse subset of reliable velocity vectors or create structured missing regions.
    \item \textbf{Thermocouples provide direct point temperature measurements, but cannot resolve 2D instantaneous fields:} Physical thermal probes (such as B-type or R-type platinum--rhodium thermocouples) measure temperature only at isolated spatial locations ($K \sim 1\text{--}32$ probes). Inserting a dense rake of thermocouples into the reacting core disturbs the swirling vortex breakdown, introduces radiative heat-loss bias ($\propto \epsilon\sigma(T_{\text{probe}}^4 - T_{\text{wall}}^4)$) and catalytic radical recombination on bare wire surfaces, suffers from thermal inertia ($\tau_{\text{tc}} \sim 10\text{--}100\text{ ms}$) that filters out kilohertz-scale turbulent fluctuations, and risks probe failure in stoichiometric methane-rich zones ($T > 2100\text{ K}$).
    \item \textbf{Chemiluminescence ($\text{OH}^*$) is line-of-sight integrated, not planar temperature:} High-speed cameras equipped with UV intensifiers and bandpass filters around $308\text{ nm}$ routinely record emission from electronically excited hydroxyl radicals ($\text{OH}^*$, $A^2\Sigma^+ \rightarrow X^2\Pi$) at kilohertz frame rates. However, $\text{OH}^*$ is short-lived and forms only within the sub-millimeter exothermic reaction sheet before undergoing rapid collisional quenching, and its optical signal is integrated along the entire line of sight through the three-dimensional annular flame. Consequently, $\text{OH}^*$ imaging provides neither a two-dimensional planar slice nor quantitative temperatures inside the cold reactant jets or the post-flame Inner Recirculation Zone (where ground-state $\text{OH}$ persists without chemiluminescent emission).
    \item \textbf{Planar laser thermometry and 3D Large Eddy Simulations (LES) are prohibitively expensive across large test matrices:} Two-line planar laser-induced fluorescence (OH-PLIF), Rayleigh scattering, or coherent anti-Stokes Raman spectroscopy (CARS) require complex multi-laser optical chains and species-dependent collisional-quenching calibrations~\cite{kohse2005applied}, while high-fidelity 3D Large Eddy Simulations with reduced finite-rate chemistry (e.g., 38 species and 184 elementary reactions~\cite{hu2026realpdebench}) require thousands of CPU/GPU hours per operating point.
\end{enumerate}

This diagnostic asymmetry motivates a concrete inverse problem for laboratory testing: \emph{if a set of high-fidelity LES simulations (or calibrated experiments) is available on a training set of operating conditions, can a data-driven model learn to reconstruct the hidden two-dimensional instantaneous temperature field $T(\mathbf{x}, t)$ at new time instants---and at unseen fuel blends $(\text{NH}_3\%, \phi)$---using only sparse velocity observations, optionally aided by a few point temperature sensors?}

\section{Problem Formulation and Research Questions}
\label{sec:intro_questions}

At first glance, inferring temperature from velocity alone may appear physically ill-posed. In a low-Mach reacting flow, velocity and temperature are coupled through Continuity (thermal gas expansion across the flame front, $\nabla \cdot \mathbf{u} = -\frac{1}{\rho}\frac{D\rho}{Dt}$), buoyancy, and the geometry of the swirl-induced shear layers: high-speed inflowing jets correspond to cold unburnt reactants ($T \approx 300\text{ K}$), whereas low-speed recirculation pockets contain hot burnt products. However, this kinematic-thermal coupling suffers from two fundamental sources of ambiguity:
\begin{enumerate}
    \item \textbf{Cross-condition thermodynamic ambiguity:} Across different $\text{NH}_3/\text{CH}_4$ fuel blends and equivalence ratios $\phi$, the bulk inflow velocity and swirl number remain similar ($\sim 10\text{--}15\text{ m/s}$), whereas the adiabatic flame temperature and burnt-gas temperature level shift by $200\text{--}360\text{ K}$. Consequently, an identical local velocity vector can correspond to $2150\text{ K}$ in pure methane or $1800\text{ K}$ in an $80\%\text{ NH}_3$ blend.
    \item \textbf{Instantaneous turbulent flame-brush ambiguity:} While time-averaged temperature fields vary smoothly with $(\text{NH}_3\%, \phi)$, instantaneous turbulent snapshots exhibit intense vortex shedding and flame-front wrinkling ($\sigma_T \approx 300\text{--}500\text{ K}$ along the jet shear layers). Linear modal estimators such as Gappy Proper Orthogonal Decomposition (Gappy POD)~\cite{everson1995karhunen} capture the large-scale mean structure but struggle to infer the nonlinear phase and position of the wrinkled reaction front from sparse velocity vectors.
\end{enumerate}

To study both sources of ambiguity systematically without conflating them, this thesis separates the reconstruction problem into two distinct operational tracks:
\begin{itemize}
    \item \textbf{Track 1 --- Temperature-Free Reconstruction ($\mathcal{O}_u \rightarrow T$):} At inference time, no temperature sensors are available ($K = 0$). The model observes only sparse velocity vectors $\mathcal{O}_u = \{(\mathbf{x}_m, \mathbf{u}(\mathbf{x}_m, t))\}_{m=1}^{M_u}$ on the 2D plane, either without operating metadata (\emph{unconditioned}) or conditioned on the known burner operating parameters $\mathbf{p} = (\text{NH}_3\%, \phi)$ (\emph{parameter-conditioned}).
    \item \textbf{Track 2 --- Temperature-Assisted Reconstruction ($\mathcal{O}_u \cup \mathcal{O}_T \rightarrow T$):} In addition to sparse velocity $\mathcal{O}_u$ (and optional $\mathbf{p}$), the model observes $K \in \{1, 4, 9, 16, 32\}$ sparse point temperature measurements $\mathcal{O}_T = \{(\mathbf{x}_k, T(\mathbf{x}_k, t))\}_{k=1}^K$, mimicking a small array of physical thermocouples.
\end{itemize}

Within this framework, the thesis addresses four specific scientific questions:
\begin{enumerate}
    \item \textbf{Mean Thermal Level vs.\ Instantaneous Flame-Front Structure:} When reconstructors are evaluated across seen conditions, held-out interpolation conditions (\texttt{test\_interp}), and held-out extrapolation conditions at the lean/rich and high-$\text{NH}_3$ boundaries (\texttt{test\_extrap}), how much of the reconstruction error stems from a uniform shift in global thermal level ($\text{Bias}$) versus spatial pattern error along the turbulent flame brush ($\text{Centered-RMSE}$ and regional jet/front error)?
    \item \textbf{Representation Ceiling vs.\ Linear Estimation Gap:} Can a low-dimensional linear basis (POD) represent the instantaneous temperature fluctuations across operating conditions if the true temperature field is known (\texttt{OraclePOD}), and how large is the estimation gap when those modal coefficients are inferred linearly from sparse velocity (\texttt{GappyPOD}) versus nonlinearly via deep neural architectures?
    \item \textbf{Global-Local Neural Reconstruction under Sparsity:} How effectively can a coordinate-conditioned global-local token architecture---adapted from the laboratory's \emph{Global-Local-Uncertainty} (\textbf{GLU}) framework~\cite{wang2026glu, wang2026ammonia, wang2026dmfgen}---combine global operating context with local neighborhood velocity features to reconstruct wrinkled instantaneous flame fronts across sparsity levels ($1\%\text{--}20\%$), and how does it compare against convolutional masked-grid baselines (U-Net) and classical modal methods?
    \item \textbf{Assimilation of Small Thermocouple Arrays ($K > 0$):} How many sparse temperature anchors $K$ are required to calibrate global thermal bias versus local flame-front shape, and how can small thermocouple counts ($K \le 9$) be assimilated into a velocity-conditioned neural reconstructor without being overwhelmed by orders-of-magnitude more velocity tokens ($M_u \gg K$)?
\end{enumerate}

\section{Methodological Guardrails and Thesis Outline}
\label{sec:intro_outline}

To ensure that all quantitative conclusions reflect genuine physical generalization rather than data leakage or numerical artifacts, the experimental pipeline enforces four strict guardrails throughout the thesis:
\begin{enumerate}
    \item \textbf{Condition-level splitting and temporal buffering:} Because local temperature fluctuations in the flame brush remain temporally correlated over $50\text{--}100$ frames ($\tau_c \approx 14\text{ ms}$ at $4000\text{ Hz}$), entire operating-condition trajectories are withheld into disjoint training (\texttt{train}, 20 conditions), validation (\texttt{val}, 3 conditions), interpolation (\texttt{test\_interp}, 3 conditions), and extrapolation (\texttt{test\_extrap}, 4 conditions) splits. Seen-condition diagnostics on training trajectories are evaluated on a future time block separated from training frames by a buffer of at least $120$ frames ($\ge 2\tau_c$).
    \item \textbf{Strict fit-only statistics and input isolation:} All normalization statistics, parametric mean priors, and POD bases are fitted exclusively on the designated training frames (\texttt{train\_fit}). At inference time, the target temperature field $T$ is strictly isolated from model inputs.
    \item \textbf{Explicit sparse representations:} Unobserved grid locations are never passed to models as physical zeros; sparse observations are represented either as coordinate-value tokens with validity masks or as explicit value-plus-mask channels.
    \item \textbf{Physical channel fidelity:} Planar LES hydroxyl mole fraction ($\text{OH}$, Ch~8) is never equated with line-of-sight integrated $\text{OH}^*$ chemiluminescence, and Euclidean velocity magnitude $|\mathbf{u}|$ is always recomputed directly from the orthogonal velocity components $(v_i, v_j, v_k)$ to avoid dataset interpolation artifacts.
\end{enumerate}

The remainder of this thesis is organized into five chapters:
\begin{itemize}
    \item \textbf{Chapter~\ref{ch:background} (Physical and Methodological Background)} reviews the fluid mechanics and thermochemistry of swirl-stabilized $\text{NH}_3/\text{CH}_4$ flames, classical sparse field reconstruction methods (POD, Gappy POD, and local regression), physics-informed neural networks for reacting flows (\textit{RF-PINNs}~\cite{yadav2025rfpinns}) and why naive 2D PDE residual losses are ill-suited to instantaneous 3D swirling slices, and the mathematical formulation of the \textit{GLU} architecture~\cite{wang2026glu}.
    \item \textbf{Chapter~\ref{ch:dataset} (Dataset Audit, Physical Coupling Analysis, and Experimental Protocol)} details the 30 LES operating conditions of the \textit{RealPDEBench} combustion benchmark~\cite{hu2026realpdebench}, establishes the geometric and channel conventions of the 2D diametral slice, quantifies the kinematic-thermal coupling and cross-condition ambiguity across blends, and defines the locked validation splits, observation masks, and decomposed evaluation metrics.
    \item \textbf{Chapter~\ref{ch:methodology} (Reconstruction Methodology: Classical Baselines and Neural Architectures)} presents the mathematical formulation of all reconstruction models evaluated in this work, progressing from non-neural baselines (\texttt{GlobalTrainMean}, the LOCO-selected \texttt{ParamConditionedMean}, \texttt{LocalHGB}, variance-scaled \texttt{GappyPOD}, \texttt{AnchorOffset}, and the \texttt{OraclePOD} linear representation ceiling) to dense and masked U-Net architectures and our adaptations of the \textit{GLU} token-based reconstructor for temperature-free and temperature-assisted inference.
    \item \textbf{Chapter~\ref{ch:results} (Experimental Results and Discussion)} reports the quantitative and qualitative comparisons across all splits, beginning with the classical baseline benchmarks and linear representation bounds, followed by dense velocity diagnostics, sparse temperature-free neural reconstruction across coverage levels, and temperature-assisted assimilation across nested thermocouple counts $K$.
    \item \textbf{Chapter~\ref{ch:conclusions} (Conclusions and Future Perspectives)} summarizes the main scientific findings, practical implications for laboratory combustion diagnostics, and directions for future extensions including temporal forecasting and optimal sensor placement.
\end{itemize}
